% Generated draft: requires mathematical and source review.
\documentclass[11pt]{article}
\usepackage[margin=25mm]{geometry}
\usepackage{fontspec}
\setmainfont{Latin Modern Roman}
\usepackage{amsmath,amssymb,mathtools,unicode-math}
\setmathfont{Latin Modern Math}
\usepackage{enumitem,needspace,xurl,hyperref}
\setlength{\parindent}{0pt}
\setlength{\parskip}{5pt}
\setlength{\emergencystretch}{3em}
\newcommand{\N}{\mathbb{N}}
\newcommand{\Z}{\mathbb{Z}}
\newcommand{\Q}{\mathbb{Q}}
\newcommand{\R}{\mathbb{R}}
\newcommand{\C}{\mathbb{C}}
\newcommand{\F}{\mathbb{F}}
\newcommand{\A}{\mathbb{A}}
\newcommand{\PP}{\mathbb{P}}
\newcommand{\E}{\mathbb{E}}
\newcommand{\eps}{\varepsilon}
\newcommand{\dd}{\,\mathrm{d}}
\begin{document}
\section*{UnsolvedMath ID: 30003742}
\textbf{Source problem code:} catalogue entry
\subsection{Definitions and mathematical statement}
The supplied record asks about uniqueness of homological primes over triangular primes, but does not specify size or rigidity hypotheses on the tensor-triangulated category. The following is a precise formulation in the conventional essentially small setting; its intended scope requires confirmation against the source.

Let \(\mathcal K\) be an essentially small tensor-triangulated category: a triangulated category with a symmetric monoidal product \(\otimes\), unit \(\mathbf 1\), and tensor product exact in each variable. A triangular prime is a proper thick tensor ideal \(P\subsetneq\mathcal K\) such that
\[
x\otimes y\in P\quad\Longrightarrow\quad x\in P\ \text{or}\ y\in P.
\]
Here thick means closed under isomorphisms, shifts, distinguished triangles, and direct summands, and being a tensor ideal means that \(x\in P\) implies \(x\otimes y\in P\) for every \(y\in\mathcal K\).

Write \(\operatorname{mod}\mathcal K\) for the category of finitely presented additive functors \(\mathcal K^{\mathrm{op}}\to\mathrm{Ab}\). Thus an object \(M\) admits a presentation
\[
\mathcal K(-,a)\longrightarrow\mathcal K(-,b)\longrightarrow M\longrightarrow0
\]
for objects \(a,b\in\mathcal K\). Let
\[
h\colon\mathcal K\longrightarrow\operatorname{mod}\mathcal K,
\qquad h(x)=\mathcal K(-,x)
\]
be the Yoneda functor. Equip \(\operatorname{mod}\mathcal K\) with the induced right-exact tensor product, characterized on representable functors by \(h(x)\otimes h(y)\cong h(x\otimes y)\).

A Serre tensor ideal \(\mathcal B\) is a full subcategory of \(\operatorname{mod}\mathcal K\), closed under isomorphisms, subobjects, quotients, and extensions, and satisfying \(M\otimes N\in\mathcal B\) whenever \(M\in\mathcal B\) and \(N\in\operatorname{mod}\mathcal K\). In this formulation, a homological prime means a proper Serre tensor ideal maximal among all proper Serre tensor ideals. Define
\[
h^{-1}(\mathcal B)=\{x\in\mathcal K\mid h(x)\in\mathcal B\}.
\]

For each triangular prime \(P\), set
\[
\mathscr H_P=\{\mathcal B\subsetneq\operatorname{mod}\mathcal K\mid
\mathcal B\text{ is a maximal proper Serre tensor ideal and }h^{-1}(\mathcal B)=P\}.
\]
\textbf{Uniqueness question.} For every category \(\mathcal K\) in the intended class and every triangular prime \(P\subsetneq\mathcal K\), must
\[
\mathcal B_1,\mathcal B_2\in\mathscr H_P
\quad\Longrightarrow\quad
\mathcal B_1=\mathcal B_2?
\]
The supplied wording ``is there a unique'' can also be read as asking \(\lvert\mathscr H_P\rvert=1\), which additionally requires existence. The record does not explicitly distinguish that reading from uniqueness conditional on existence. It also does not specify whether every object of \(\mathcal K\) must be rigid, meaning dualizable. These qualifications remain unresolved rather than being asserted as source hypotheses.
\subsection{Short English statement}
Given a triangular prime \(P\), can two distinct maximal proper Serre tensor ideals have the same Yoneda preimage \(P\)? The record asks whether this is impossible. The intended category hypotheses and whether existence is included require clarification. The supplied status is partially solved; this draft does not independently verify that assessment.
\subsection{Sources}
\begin{itemize}
\item Local registry: \texttt{data/unsolvedmath/problems.json}, with \texttt{catalog/problems.json} as the additional-record registry; permanent UnsolvedMath ID \(30003742\), source problem code \texttt{catalogue entry}, title ``Uniqueness of Homological Primes over Triangular Primes.'' These are registry references supplied by the project instructions; no independent inspection is claimed.
\item Supplied source citation: \emph{Mini-Workshop: Chromatic Phenomena and Duality in Homotopy Theory and Representation Theory} (2019), Oberwolfach Reports, pp.~507--529, DOI \texttt{10.4171/owr/2018/9}. Supplied source URL: \url{https://doi.org/10.4171/owr/2018/9}. The record gives proposed year 2018 and names no individual proposer. The relationship between that year and the citation's 2019 date is unresolved.
\item Supplied literature source, ``original OWR report'': \url{https://doi.org/10.4171/OWR/2018/9}.
\item Supplied literature source, ``Balmer--Krause--Stevenson, tensor-triangular fields'': \url{https://doi.org/10.1007/s00029-019-0454-2}.
\item Supplied literature source, ``Balmer--Cameron, computed homological residue fields'': \url{https://arxiv.org/abs/2007.04485}.
\item Literature assessment dated exactly \texttt{2026-08-21}: supplied dated curation, not independently verified. It labels the problem partially solved and reports uniqueness in many standard stable-homotopy, algebraic-geometry, and representation-theoretic examples, while reporting no known general uniqueness theorem. The background describes this as literature triage derived from the supplied OWR curation and requests independent verification of time-sensitive status and cited claims.
\item The registry labels the statement \texttt{corrected\_verified} and says mathematical content was preserved while notation, typography, and sentence structure were repaired for publication. This is supplied editorial metadata, not independent mathematical verification of this draft.
\end{itemize}
\end{document}
